
%% ============================================================
%% REVISED ABSTRACT
%% ============================================================
\begin{abstract}
\noindent I examine whether ideological distance between a CEO and the
firm's employees affects corporate restructuring decisions. Using campaign
contribution records from the Database on Ideology, Money in Politics, and
Elections (DIME) to measure the political ideology of both CEOs and
employees, I construct a firm--election-cycle panel of CEO--employee
political misalignment for over 1,000 publicly traded firms from 2000 to
2024. I find that misalignment constrains restructuring activity during
normal industry conditions: a one-standard-deviation increase in
misalignment is associated with a 0.5 percentage point reduction in
restructuring charges relative to assets when the firm's industry is not
experiencing widespread layoffs. This constraint is robust to industry
$\times$ year fixed effects, survives first-differences estimation using
only within-firm changes in misalignment, and exhibits clean pre-trends
with no evidence of reverse causality. The constraint is concentrated in
periods of low industry peer layoff activity and attenuates when
industry-wide distress is high, consistent with peer behavior providing
external justification that overcomes organizational resistance. The effect
is symmetric regardless of whether the CEO is to the right or left of
employees (Wald test $p = 0.92$), confirming that ideological
\textit{distance}---not direction---generates friction. Cumulative two-year
charges are unaffected by misalignment, indicating the effect operates on
the \textit{timing} of restructuring rather than the total level. These
findings suggest that political misalignment functions as organizational
friction that delays but does not permanently prevent restructuring.
\medskip
\noindent \textit{JEL Codes:} G34, M41, D72, J50 \\
\noindent \textit{Keywords:} Political ideology, CEO--employee misalignment,
restructuring, campaign contributions, DIME
\end{abstract}

%% ============================================================
%% REVISED INTRODUCTION PARAGRAPH (replace para beginning
%% "The main findings are as follows")
%% ============================================================
The main findings are as follows. First, the unconditional relationship
between misalignment and restructuring is negative but not statistically
significant with firm and year fixed effects (coefficient $= -0.002$,
$p = 0.20$). Second, this average null masks meaningful heterogeneity
by industry conditions. When the firm's industry has low peer layoff
activity, misalignment significantly constrains restructuring
($-0.0094$, $p = 0.001$); this constraint attenuates and becomes
statistically indistinguishable from zero as industry-wide layoff rates
rise. Third, the constraint survives the inclusion of industry $\times$
year fixed effects: at the 10th percentile of industry peer layoff rates,
the net effect of misalignment is $-0.0047$ ($p = 0.036$), and the
pattern of net effects across quantiles is consistent with the baseline
specification. Fourth, first-differences estimation---which uses only
within-firm changes in misalignment and eliminates all time-invariant
characteristics by construction---produces a directionally consistent
coefficient ($-0.003$, $p \approx 0.14$), and pre-trend tests confirm
the effect does not precede the theorized timing. Fifth, the effect is
perfectly symmetric: it operates identically regardless of whether the
CEO is to the right or left of employees (Wald test $p = 0.92$),
confirming that ideological \textit{distance}, not direction, drives
the result. Sixth, cumulative two-year restructuring charges are
unaffected by misalignment, indicating the friction affects the
\textit{timing} of restructuring---delaying charges during normal
industry conditions---but not the total level.

These results are consistent with a ``constrained-in-normal-times''
interpretation. Ideological distance between the CEO and workforce
creates organizational friction---reduced trust, impaired communication,
and heightened resistance to top-down directives perceived as coming
from an ideological out-group. This friction delays restructuring during
stable industry conditions. When industry-wide distress is elevated,
this constraint attenuates, consistent with external conditions reducing
the perceived political content of restructuring decisions and making
employee resistance less tenable. We note that the interaction between
misalignment and industry peer layoff rates, while significant in the
baseline specification, is not separately identified under industry
$\times$ year fixed effects (since the peer layoff rate varies at the
industry--year level and is absorbed by the fixed effects). We therefore
treat the moderation result as mechanism-consistent evidence rather than
an independently identified causal claim.

%% ============================================================
%% REVISED RESULTS SECTION 5.2
%% (replace current "Peer Layoff Moderation" subsection)
%% ============================================================
\subsection{Industry Conditions as a Moderator}

Table~\ref{tab:peer} presents the central finding. Column~(1) reports
the baseline peer layoff interaction: the coefficient on misalignment
is $-0.0094$ ($p = 0.001$), indicating that misalignment significantly
constrains restructuring when peer layoff activity is low. The net
effect of misalignment at the 10th percentile of peer layoff rates is
$-0.0057$ ($p = 0.004$); at the median it is $-0.0027$ ($p = 0.16$);
and at the 90th percentile the effect is statistically indistinguishable
from zero. The crossover point---where the marginal effect of
misalignment equals zero---occurs at a peer layoff rate of 0.329,
corresponding to the 81st percentile.

The result is robust to controlling for employee ideology (Column~2),
CEO turnover (Column~3), and double-clustering by firm and year
(Column~4: interaction $p = 0.043$).

Column~(5) includes industry $\times$ year fixed effects, which absorb
the peer layoff rate main effect entirely (since this variable varies
at the industry--year level). Under this more demanding specification,
the main effect of misalignment is $-0.0065$ ($p = 0.08$), and the
net effects at low peer layoff rates remain consistent with the baseline:
$-0.0047$ ($p = 0.036$) at the 10th percentile and $-0.0032$
($p = 0.051$) at the median. The interaction coefficient under
industry $\times$ year FE should be interpreted with caution---the
peer layoff moderator is an industry-year variable, so the interaction
captures residual within-industry-year variation after absorbing the
fixed effects rather than the full cross-industry-year comparison in
Columns~(1)--(4). The consistency of the net effects across
specifications supports the robustness of the constraint mechanism.

[Table~\ref{tab:peer} about here.]

[Figure~\ref{fig:marginal} about here.]

%% ============================================================
%% REVISED RESULTS SECTION 5.3
%% (replace current "Real Outcomes: Layoff Indicator" subsection)
%% ============================================================
\subsection{Real Outcomes: Employment Changes}

The layoff indicator does not exhibit the same timing pattern as
restructuring charges. Table~\ref{tab:layoff_timing} in the Appendix
reports the lead/lag structure: the only marginally significant
coefficient ($p = 0.097$) appears at $t-1$---before misalignment is
measured---with the opposite sign from the hypothesized direction.
The main lagged specification ($t+1$) yields a null coefficient.
This timing pattern is inconsistent with the organizational friction
mechanism and suggests that layoff decisions respond to contemporaneous
conditions rather than lagged misalignment. We accordingly treat the
employment outcomes as supplementary and focus the primary analysis
on restructuring charges.

%% ============================================================
%% REVISED ROBUSTNESS SECTION: ADD FIRST DIFFERENCES
%% (add as new subsection 6.4, before Mechanism Tests)
%% ============================================================
\subsection{First Differences}

Table~\ref{tab:fd} estimates first-differenced specifications that
eliminate all time-invariant firm characteristics by construction,
using only within-firm \textit{changes} in misalignment for
identification. The coefficient on $\Delta|$Misalignment$|$ is
$-0.0031$ ($p \approx 0.14$), directionally consistent with the
firm fixed effects estimates ($-0.002$ to $-0.004$ across
specifications) and with the first-differences result from the
pre-trend analysis ($-0.0049$, $p = 0.076$). While not significant
at conventional levels, the first-differences estimate is reassuring:
it rules out the concern that the firm fixed effects result is driven
by time-invariant cross-sectional differences between high- and
low-misalignment firms rather than within-firm changes in
misalignment over time.

%% ============================================================
%% REVISED ROBUSTNESS SECTION: PRE-TRENDS
%% (add as new subsection 6.5)
%% ============================================================
\subsection{Pre-Trend Tests}

Table~\ref{tab:leadlag} reports the coefficient on misalignment from
separate regressions where the dependent variable is restructuring
charges at $t-1$, $t$, $t+1$, and $t+2$ relative to the misalignment
measure. The placebo coefficients at $t-1$ and $t$ are null and
incorrectly signed (positive), while the main effect appears only at
the theorized lag ($t+1$, coefficient $= -0.0019$). This pattern
confirms that the result is not driven by reverse causality or
anticipatory effects: restructuring does not predict subsequent
misalignment, and misalignment does not predict restructuring before
the hypothesized timing.

The interaction pre-trend test (Table~\ref{tab:placebo}) further
confirms the result is not an artifact of the industry--year correlation
structure. The true interaction coefficient of $+0.0285$ is 5.7
standard deviations above the mean of the placebo distribution
constructed from 500 random permutations of industry--year peer
layoff rates ($p < 0.002$).

